\documentclass[12pt,a4paper]{amsart}
\usepackage{amsmath,amssymb,amsthm}
\usepackage{mathtools}
\usepackage[utf8]{inputenc}
\usepackage[T1]{fontenc}
\usepackage[ngerman]{babel}
\usepackage{hyperref}
\usepackage{enumitem,booktabs,array}
\usepackage{geometry}
\geometry{margin=2.5cm}

\newtheorem{theorem}{Satz}[section]
\newtheorem{lemma}[theorem]{Lemma}
\newtheorem{korollar}[theorem]{Korollar}
\newtheorem{proposition}[theorem]{Proposition}
\newtheorem{conjecture}[theorem]{Vermutung}
\theoremstyle{definition}
\newtheorem{definition}[theorem]{Definition}
\newtheorem{beispiel}[theorem]{Beispiel}
\newtheorem{bemerkung}[theorem]{Bemerkung}

\title{Transzendenztheorie:\\
Klassische Resultate, offene Probleme und algorithmische Aspekte}
\author{Michael Fuhrmann}
\address{Specialist Mathematics Project, Build 141}
\email{specialist-maths@project.local}
\date{12. März 2026}
\keywords{transzendente Zahlen, Liouvillescher Satz, Hermite, Lindemann,
Gelfond-Schneider, Bakerscher Satz, Euler-Mascheroni-Konstante,
Schanuels Vermutung}
\subjclass[2020]{11J81, 11J91, 11J04, 11J82, 11J89}

\begin{abstract}
Die Transzendenztheorie untersucht, welche komplexen Zahlen keine Nullstellen
eines von null verschiedenen Polynoms mit rationalen Koeffizienten sind.
Dieser Artikel gibt einen Überblick über die klassischen Meilensteine:
Liouvilles Konstruktion der ersten beweisbar transzendenten Zahlen (1844),
Hermites Beweis der Transzendenz von $e$ (1873), Lindemanns Beweis der
Transzendenz von $\pi$ (1882) -- mit der Folgerung, dass die Quadratur des
Kreises unmöglich ist -- den Satz von Lindemann--Weierstrass, den Satz von
Gelfond--Schneider, der Hilberts siebtes Problem löst (1934/1935), sowie
Bakers weitreichende Verallgemeinerung auf Linearkombinationen von Logarithmen
(1966). Daneben wird die Euler-Mascheroni-Konstante $\gamma$ behandelt,
deren Irrationalität eines der bekanntesten offenen Probleme der Mathematik
darstellt, zusammen mit weiteren offenen Fragen um Schanuels Vermutung sowie
die Transzendenz von $e + \pi$ und $e \cdot \pi$. Algorithmische Aspekte
werden anhand der Python-Module \texttt{transcendence\_theory.py} und
\texttt{euler\_gamma\_irrationality.py} illustriert.
\end{abstract}

\maketitle

\tableofcontents

%% ==========================================================================
\section{Einleitung}
%% ==========================================================================

Die Frage, welche Zahlen als Nullstellen von Polynomgleichungen mit rationalen
Koeffizienten dargestellt werden können, gehört zu den ältesten und tiefsten
der Mathematik. Eine komplexe Zahl $\alpha$ heißt \emph{algebraisch} (über
$\mathbb{Q}$), wenn ein von null verschiedenes Polynom $p(x) \in \mathbb{Z}[x]$
existiert mit $p(\alpha) = 0$; andernfalls heißt $\alpha$ \emph{transzendent}.

Die Geschichte der Transzendenztheorie umspannt fast zwei Jahrhunderte:
\begin{itemize}[leftmargin=2em]
  \item \textbf{1844}: Joseph Liouville zeigte als Erster die Existenz
        transzendenter Zahlen, indem er explizite Beispiele konstruierte ---
        die Liouville-Zahlen ---, deren Approximierbarkeit durch rationale
        Zahlen zu schnell ist, um algebraisch zu sein.
  \item \textbf{1851}: Liouville veröffentlichte sein allgemeines Kriterium
        (Satz von Liouville über diophantische Approximation).
  \item \textbf{1873}: Charles Hermite bewies, dass die Basis des natürlichen
        Logarithmus, $e \approx 2{,}71828\ldots$, transzendent ist --- der
        erste Beweis, dass eine \emph{spezifisch auftretende} Konstante
        transzendent ist.
  \item \textbf{1882}: Ferdinand von Lindemann bewies die Transzendenz von
        $\pi \approx 3{,}14159\ldots$ und löste damit die jahrtausendealte
        Frage nach der Quadratur des Kreises: Es ist unmöglich, mit Zirkel
        und Lineal ein Quadrat mit dem Flächeninhalt $\pi$ zu konstruieren.
  \item \textbf{1885}: Karl Weierstrass verallgemeinerte Lindemanns Methode
        zum Satz von Lindemann--Weierstrass.
  \item \textbf{1900}: David Hilbert formulierte die Transzendenz von
        $2^{\sqrt{2}}$ als siebtes Problem seiner berühmten Liste von 23
        offenen Problemen.
  \item \textbf{1934/1935}: Alexander Gelfond und Theodor Schneider lösten
        unabhängig voneinander Hilberts siebtes Problem.
  \item \textbf{1966}: Alan Baker bewies eine weitreichende Verallgemeinerung
        über lineare Formen in Logarithmen algebraischer Zahlen und erhielt
        dafür 1970 die Fields-Medaille.
\end{itemize}

Trotz dieser großen Fortschritte bleiben viele grundlegende Fragen offen.
Ob $\gamma = 0{,}5772\ldots$ (die Euler-Mascheroni-Konstante) überhaupt
irrational ist, ist unbekannt. Die Transzendenz von $e + \pi$, $e \cdot \pi$,
$\zeta(3)/\pi^3$ und vieler anderer natürlich auftretender Konstanten harrt
des Beweises.

Dieser Artikel ist wie folgt gegliedert: Abschnitt~\ref{sec:alg-transz}
erläutert die Definitionen und Abzählbarkeitsargumente.
Abschnitte~\ref{sec:liouville} bis~\ref{sec:baker} stellen die wichtigsten
Sätze in historischer Reihenfolge vor. Abschnitt~\ref{sec:offen} sammelt die
offenen Probleme, und Abschnitt~\ref{sec:euler-mascheroni} widmet sich der
Euler-Mascheroni-Konstante.


%% ==========================================================================
\section{Algebraische und transzendente Zahlen}
\label{sec:alg-transz}
%% ==========================================================================

\begin{definition}
Eine komplexe Zahl $\alpha \in \mathbb{C}$ heißt \emph{algebraisch} (über
$\mathbb{Q}$), falls ein von null verschiedenes Polynom $p(x) = a_n x^n +
\cdots + a_0 \in \mathbb{Z}[x]$ mit $p(\alpha) = 0$ existiert. Der
\emph{algebraische Grad} von $\alpha$ ist der Grad ihres Minimalpolynoms,
d.\,h.\ des normierten Polynoms kleinsten Grades in $\mathbb{Q}[x]$, das bei
$\alpha$ verschwindet.

Eine nicht-algebraische Zahl heißt \emph{transzendent}.
\end{definition}

\begin{beispiel}
\begin{enumerate}[label=(\roman*)]
  \item Jede rationale Zahl $r = p/q$ ist algebraisch vom Grad 1.
  \item $\sqrt{2}$ ist algebraisch vom Grad 2, Minimalpolynom $x^2 - 2$.
  \item Jede $n$-te Einheitswurzel $\zeta_n = e^{2\pi i/n}$ ist algebraisch
        (Nullstelle des Kreisteilungspolynoms $\Phi_n(x)$).
  \item $\sqrt{2} + \sqrt{3}$ ist algebraisch vom Grad 4.
  \item $e$, $\pi$, $e^\pi$ sind transzendent (s.\ spätere Abschnitte).
\end{enumerate}
\end{beispiel}

Die Menge $\overline{\mathbb{Q}}$ aller algebraischen Zahlen bildet einen
algebraisch abgeschlossenen Teilkörper von $\mathbb{C}$. Ihre Mächtigkeit
ist abzählbar.

\begin{proposition}
Die Menge der algebraischen Zahlen $\overline{\mathbb{Q}}$ ist abzählbar.
\end{proposition}
\begin{proof}
Für jeden Grad $d \geq 1$ und jede Koeffizientenwahl $(a_0, \ldots, a_d) \in
\mathbb{Z}^{d+1}$ mit $a_d \neq 0$ hat das Polynom $a_d x^d + \cdots + a_0$
höchstens $d$ komplexe Nullstellen. Die Menge aller solchen Polynome ist
abzählbar (als abzählbare Vereinigung abzählbarer Mengen), und jedes Polynom
liefert endlich viele algebraische Zahlen. Somit ist $\overline{\mathbb{Q}}$
eine abzählbare Vereinigung endlicher Mengen, also abzählbar.
\end{proof}

\begin{korollar}
Transzendente Zahlen existieren und bilden tatsächlich eine überabzählbare
Menge der Mächtigkeit des Kontinuums. Fast jede reelle Zahl (im Sinne des
Lebesgue-Maßes) ist transzendent.
\end{korollar}

\begin{bemerkung}
Obwohl die Existenz transzendenter Zahlen bereits aus dem Abzählbarkeitsargument
(Cantor, 1874) folgt, liefert ein solcher Beweis kein konkretes Beispiel.
Dies motivierte Liouvilles explizite Konstruktion.
\end{bemerkung}


%% ==========================================================================
\section{Liouvilles Satz und Liouville-Zahlen}
\label{sec:liouville}
%% ==========================================================================

Die ersten expliziten transzendenten Zahlen wurden von Liouville 1844 mittels
seines \emph{Satzes über diophantische Approximation} konstruiert.

\begin{theorem}[Liouville, 1844]
\label{satz:liouville}
Sei $\alpha$ eine reelle algebraische Zahl vom Grad $d \geq 1$. Dann existiert
eine Konstante $c = c(\alpha) > 0$, sodass für alle Brüche $p/q \in \mathbb{Q}$
(in gekürzter Form, $q \geq 1$) gilt:
\[
  \left|\alpha - \frac{p}{q}\right| \geq \frac{c}{q^d}.
\]
Algebraische Zahlen vom Grad $d$ können also nicht schneller als mit der Rate
$q^{-d}$ durch rationale Zahlen approximiert werden.
\end{theorem}

\begin{proof}[Beweissskizze]
Sei $f(x) = a_d x^d + \cdots + a_0 \in \mathbb{Z}[x]$ das Minimalpolynom von
$\alpha$. Da $f(\alpha) = 0$ und $f$ über $\mathbb{Q}$ irreduzibel ist, gilt
$f(p/q) \neq 0$ für alle $p/q \in \mathbb{Q}$.
Nenner-Freistellung: $|q^d f(p/q)| \geq 1$ (ganze Zahl, $\neq 0$).
Nach dem Mittelwertsatz gilt $f(p/q) - f(\alpha) = f'(\xi)(p/q - \alpha)$
für ein $\xi$ zwischen $p/q$ und $\alpha$. Für $|p/q - \alpha| \leq 1$
ist $|f'(\xi)| \leq M$ beschränkt. Daraus folgt
$|p/q - \alpha| \geq 1/(M q^d)$.
\end{proof}

Dieser Satz ermöglicht die explizite Konstruktion transzendenter Zahlen:

\begin{definition}
Eine reelle Zahl $\xi$ heißt \emph{Liouville-Zahl}, falls für jedes $n \in
\mathbb{N}$ unendlich viele Brüche $p/q$ mit $q > 1$ existieren, sodass
\[
  \left|\xi - \frac{p}{q}\right| < \frac{1}{q^n}.
\]
Liouville-Zahlen besitzen das \emph{Irrationalitätsmaß} $\mu(\xi) = \infty$.
\end{definition}

\begin{beispiel}[Liouvillesche Konstante]
Die Zahl
\[
  L = \sum_{k=1}^{\infty} 10^{-k!}
  = 0{,}1100010000000000000000010\ldots
\]
ist eine Liouville-Zahl und damit transzendent. Die Partialsummen
$p_n/q_n$ mit $q_n = 10^{n!}$ erfüllen $|L - p_n/q_n| < 2 \cdot 10^{-(n+1)!}
= 2/q_n^{n+1}$, was Liouvilles Schranke für jeden endlichen Grad $d$ verletzt.
\end{beispiel}

\begin{bemerkung}
Nicht alle transzendenten Zahlen sind Liouville-Zahlen. Zum Beispiel besitzt
$e$ das Irrationalitätsmaß $\mu(e) = 2$, ebenso wie jede algebraisch irrationale
Zahl. Die derzeit beste bekannte obere Schranke für das Irrationalitätsmaß von $\pi$ ist $\mu(\pi) \leq 7{,}606$
(Salikhov, 2008).
\end{bemerkung}

Das \emph{Irrationalitätsmaß} $\mu(x)$ ist definiert als
\[
  \mu(x) = \sup\!\left\{\mu \in \mathbb{R} : \left|x - \frac{p}{q}\right|
  < \frac{1}{q^\mu} \text{ hat unendlich viele Lösungen } p/q \in \mathbb{Q}
  \right\}.
\]
Der Satz von Dirichlet impliziert $\mu(x) \geq 2$ für alle irrationalen $x$.


%% ==========================================================================
\section{\texorpdfstring{$e$ ist transzendent}{e ist transzendent}}
\label{sec:e-transz}
%% ==========================================================================

\begin{theorem}[Hermite, 1873]
\label{satz:e-transz}
Die Zahl $e = \sum_{n=0}^{\infty} 1/n!$ ist transzendent.
\end{theorem}

\begin{proof}[Beweissskizze]
Annahme: $e$ ist algebraisch. Dann existieren ganze Zahlen $a_0, a_1, \ldots,
a_n$ mit $a_0, a_n \neq 0$, sodass
\[
  a_0 + a_1 e + a_2 e^2 + \cdots + a_n e^n = 0.
\]
Für eine große Primzahl $p$ (die geeignet gewählt wird) definiere das Hilfspolynom
\[
  f(x) = \frac{x^{p-1}(x-1)^p(x-2)^p \cdots (x-n)^p}{(p-1)!}.
\]
Setze $F(x) = \sum_{j=0}^{np+p-1} f^{(j)}(x)$ (Summe aller Ableitungen).
Betrachte das Integral
\[
  J(t) = e^t \int_0^t e^{-u} f(u)\,du = F(0)\,e^t - F(t).
\]
Definiere $I_k = \int_0^k e^{k-u} f(u)\,du$, $k = 0, \ldots, n$.
Dann gilt $\sum_{k=0}^n a_k I_k = F(0) \underbrace{\sum_{k=0}^n a_k e^k}_{=0}
- \sum_{k=0}^n a_k F(k)$.

Man zeigt: (1) $F(0)$ und $F(k)$ sind ganze Zahlen, die durch $p$ teilbar sind,
wobei $F(0)/p$ für große Primzahl $p$ (die $a_0 \cdot n!$ nicht teilt) nicht
durch $p$ teilbar ist -- also ist $\sum a_k F(k) \neq 0$.
(2) Andererseits gilt $|I_k| \to 0$ für $p \to \infty$ (der Faktor $(p-1)!$
im Nenner dominiert). Dies ergibt eine von null verschiedene ganze Zahl mit
beliebig kleinem Betrag -- Widerspruch.
\end{proof}

\begin{bemerkung}
Hermites ursprünglicher Beweis war erheblich rechenaufwändiger. Die hier
präsentierte klarere Version geht teilweise auf Hilbert und Hurwitz (ca.\ 1893)
zurück. Die zentrale Idee -- Konstruktion einer Funktion, deren Werte an
ganzzahligen Stellen nahezu ganze Zahlen sind -- wurde zum Vorbild aller
späteren Transzendenznachweise.
\end{bemerkung}

Numerisch lässt sich verifizieren, dass kein Polynom kleinen Grades mit kleinen
ganzzahligen Koeffizienten bei $e$ verschwindet. Das Modul
\texttt{transcendence\_theory.py} implementiert \texttt{e\_is\_transcendental\_demo()},
das alle normierten Polynome vom Grad $\leq 6$ mit Koeffizienten in
$\{-5, \ldots, 5\}$ prüft: Der Minimalwert von $|f(e)|$ über alle solchen
Polynome beträgt etwa $3{,}9 \times 10^{-3}$ (für Grad 2).


%% ==========================================================================
\section{\texorpdfstring{$\pi$ ist transzendent}{pi ist transzendent}}
\label{sec:pi-transz}
%% ==========================================================================

\begin{theorem}[Lindemann, 1882]
\label{satz:pi-transz}
Die Zahl $\pi$ ist transzendent.
\end{theorem}

\begin{proof}[Beweissskizze (via Lindemann--Weierstrass)]
Die Eulersche Formel $e^{i\pi} + 1 = 0$ ergibt $e^{i\pi} = -1$.
Falls $\pi$ algebraisch wäre, so auch $i\pi$ (da $i$ algebraisch ist als
Nullstelle von $x^2 + 1$). Der Satz von Lindemann--Weierstrass
(Satz~\ref{satz:lw}) besagt, dass $e^\alpha$ für jedes von null verschiedene
algebraische $\alpha$ transzendent ist. Aber $e^{i\pi} = -1$ ist algebraisch
-- Widerspruch. Also ist $\pi$ transzendent.
\end{proof}

\begin{korollar}[Quadratur des Kreises ist unmöglich]
Es ist unmöglich, allein mit Zirkel und Lineal ein Quadrat zu konstruieren,
dessen Fläche gleich der eines gegebenen Kreises ist.
\end{korollar}
\begin{proof}
Eine Länge ist genau dann konstruierbar, wenn sie algebraisch über $\mathbb{Q}$
ist und in einer iterierten quadratischen Erweiterung liegt. Da $\pi$ transzendent
ist, ist es nicht konstruierbar. Ein Quadrat mit Flächeninhalt $\pi r^2$ würde
die Konstruktion von $\sqrt{\pi} \cdot r$ erfordern, aber $\sqrt{\pi}$ ist
ebenfalls transzendent (wäre $\sqrt{\pi} = \alpha$ algebraisch, so auch
$\pi = \alpha^2$ -- Widerspruch).
\end{proof}

\begin{bemerkung}
Mit derselben Methode zeigt man, dass $\cos(\alpha)$, $\sin(\alpha)$,
$\tan(\alpha)$ für jedes von null verschiedene algebraische $\alpha$ transzendent
sind, und dass $\log(\beta)$ für jedes positive algebraische $\beta \neq 1$
transzendent ist.
\end{bemerkung}


%% ==========================================================================
\section{Der Satz von Lindemann--Weierstrass}
\label{sec:lw}
%% ==========================================================================

\begin{theorem}[Lindemann--Weierstrass, 1882/1885]
\label{satz:lw}
Seien $\alpha_1, \ldots, \alpha_n \in \overline{\mathbb{Q}}$ paarweise
verschiedene algebraische Zahlen. Dann sind $e^{\alpha_1}, \ldots, e^{\alpha_n}$
linear unabhängig über $\overline{\mathbb{Q}}$.

Äquivalent: Sind $\beta_1, \ldots, \beta_n \in \overline{\mathbb{Q}}$ nicht alle
null, so gilt
\[
  \beta_1 e^{\alpha_1} + \cdots + \beta_n e^{\alpha_n} \neq 0.
\]
\end{theorem}

Ein wichtiger Spezialfall:

\begin{korollar}
Für jedes von null verschiedene algebraische $\alpha$ ist die Zahl $e^\alpha$
transzendent.
\end{korollar}

Der Beweis des Satzes von Lindemann--Weierstrass verwendet eine
Hilfspolynomkonstruktion analog zu Hermites Beweis, aber für gleichzeitig
mehrere Exponentialfunktionen. Wesentlich ist eine Abschätzung der Form:
Würde die Schlussfolgerung scheitern, ergäbe sich ein ganzzahliger Ausdruck,
der von null verschieden, aber beliebig klein ist -- ein Widerspruch.

\begin{beispiel}
\begin{enumerate}[label=(\roman*)]
  \item $\alpha = 1$: $e^1 = e$ ist transzendent (Hermite).
  \item $\alpha = i\pi$: $e^{i\pi} = -1$ ist algebraisch; per Kontraposition
        ist $i\pi$ transzendent, und da $i$ algebraisch ist, ist $\pi$
        transzendent.
  \item $\alpha = \sqrt{2}$: Da $\sqrt{2}$ algebraisch und von null verschieden
        ist, ist $e^{\sqrt{2}}$ transzendent.
  \item $\ln 2$ ist transzendent: Wäre $\ln 2$ algebraisch, so ergäbe der Satz
        $e^{\ln 2} = 2$ -- aber $2$ ist algebraisch, die Kontraposition liefert
        also, dass $\ln 2$ transzendent sein muss.
\end{enumerate}
\end{beispiel}

\begin{korollar}
Für $r \in \mathbb{Q} \setminus \{0\}$ sind $\sin(r)$, $\cos(r)$, $\tan(r)$,
$e^r$, $\ln(r+1)$ alle transzendent.
\end{korollar}


%% ==========================================================================
\section{Der Satz von Gelfond--Schneider}
\label{sec:gelfond-schneider}
%% ==========================================================================

Hilberts siebtes Problem (1900) fragte: Wenn $\alpha$ und $\beta$ algebraisch
sind, $\alpha \neq 0, 1$, und $\beta$ irrational, muss dann $\alpha^\beta$
transzendent sein?

\begin{theorem}[Gelfond--Schneider, 1934/1935]
\label{satz:gelfond-schneider}
Seien $\alpha, \beta \in \overline{\mathbb{Q}}$ mit $\alpha \neq 0$, $\alpha
\neq 1$ und $\beta \notin \mathbb{Q}$. Dann ist $\alpha^\beta$ transzendent.
\end{theorem}

Alexander Gelfond (1934) und Theodor Schneider (1935) bewiesen dieses Resultat
unabhängig voneinander und lösten damit Hilberts siebtes Problem.

\begin{proof}[Beweisidee]
Man schreibt $\alpha^\beta = e^{\beta \log \alpha}$. Der Beweis konstruiert
eine Hilfsfunktion $F(z) = \sum_{j,k} c_{jk} \alpha^{jz} \beta^k z^k$
und zeigt, dass diese Funktion, wäre $\alpha^\beta$ algebraisch, zu oft
verschwinden würde -- im Widerspruch zu einer analytischen Abschätzung
(die Funktion kann nicht identisch null sein und gleichzeitig auf einem
großen Gitter verschwinden).
\end{proof}

\begin{beispiel}
\begin{enumerate}[label=(\roman*)]
  \item \textbf{Gelfond-Schneider-Konstante:} $2^{\sqrt{2}} \approx 2{,}6651\ldots$
        ist transzendent. Hier $\alpha = 2$, $\beta = \sqrt{2}$ (irrational algebraisch).
  \item \textbf{Gelfondsche Konstante:} $e^\pi = (e^{i\pi})^{-i} = (-1)^{-i}$.
        Mit $\alpha = -1$ (algebraisch, $\neq 0,1$) und $\beta = -i$
        (irrational algebraisch) ist $e^\pi$ transzendent.
        Numerisch $e^\pi \approx 23{,}1407\ldots$.
  \item $i^i = e^{i \log i} = e^{i \cdot i\pi/2} = e^{-\pi/2}$ ist transzendent
        (da $e^\pi$ transzendent ist).
\end{enumerate}
\end{beispiel}

\begin{bemerkung}
Der Satz von Gelfond--Schneider gilt nicht, wenn $\beta \in \mathbb{Q}$ (also
algebraisch vom Grad 1). Zum Beispiel ist $2^{1/2} = \sqrt{2}$ algebraisch.
Die Bedingung $\beta \notin \mathbb{Q}$ ist wesentlich.
\end{bemerkung}


%% ==========================================================================
\section{Der Bakersche Satz über lineare Formen in Logarithmen}
\label{sec:baker}
%% ==========================================================================

Bakers Satz (1966) ist eine weitreichende Verallgemeinerung des Satzes von
Gelfond--Schneider auf lineare Formen in Logarithmen algebraischer Zahlen.

\begin{theorem}[Baker, 1966]
\label{satz:baker}
Seien $\alpha_1, \ldots, \alpha_n \in \overline{\mathbb{Q}} \setminus \{0\}$
und $\beta_0, \beta_1, \ldots, \beta_n \in \overline{\mathbb{Q}}$, nicht alle
null. Gilt
\[
  \Lambda = \beta_0 + \beta_1 \log \alpha_1 + \cdots + \beta_n \log \alpha_n \neq 0,
\]
so ist $\Lambda$ transzendent.

Darüber hinaus bewies Baker eine \emph{effektive untere Schranke}:
\[
  |\Lambda| \geq \exp\!\bigl(-C \cdot (\log A_1) \cdots (\log A_n) \cdot \log B\bigr),
\]
wobei $A_i$ und $B$ geeignet definierte Höhenparameter sind und $C > 0$ nur
von $n$ und den Graden der $\alpha_i$, $\beta_j$ abhängt.
\end{theorem}

Alan Baker erhielt für diesen Satz 1970 die Fields-Medaille.

\begin{bemerkung}
Der Satz von Gelfond--Schneider entspricht dem Spezialfall $n = 1$, $\beta_0 =
0$, $\beta_1 = \beta$, $\alpha_1 = \alpha$: $\Lambda = \beta \log \alpha \neq
0$ impliziert, dass $e^\Lambda = \alpha^\beta$ transzendent ist.
\end{bemerkung}

\begin{beispiel}
\begin{enumerate}[label=(\roman*)]
  \item $\log 2$, $\log 3$ und $\log 6$ sind paarweise linear unabhängig über
        $\overline{\mathbb{Q}}$ (alle transzendent nach Lindemann--Weierstrass).
  \item Die Identität $\pi = -i \log(-1)$ zeigt mittels Baker die Transzendenz
        von $\pi$ (im komplexen Logarithmus).
  \item Bakers effektive Schranken liefern konkrete Resultate: z.\,B.\ gilt
        $|m \log 2 - n \log 3| > (mn)^{-C}$ für explizites $C$, was in der
        Theorie der $S$-Einheiten-Gleichungen entscheidend ist.
\end{enumerate}
\end{beispiel}

Bakers Satz hatte enorme Konsequenzen für die Zahlentheorie:
\begin{itemize}[leftmargin=2em]
  \item Effektive Lösung der Thue-Gleichung und der Thue-Mahler-Gleichungen.
  \item Beweis, dass jede elliptische Kurve über $\mathbb{Q}$ nur endlich viele
        ganzzahlige Punkte besitzt (effektive Version des Siegelschen Satzes).
  \item Anwendungen auf Klassenzahlen und $S$-ganzzahlige Punkte.
  \item Algorithmische Grundlage für die Lösung vieler diophantischer Gleichungen.
\end{itemize}


%% ==========================================================================
\section{Die Euler-Mascheroni-Konstante}
\label{sec:euler-mascheroni}
%% ==========================================================================

Die Euler-Mascheroni-Konstante ist definiert durch
\[
  \gamma = \lim_{n \to \infty} \left( \sum_{k=1}^n \frac{1}{k} - \ln n \right)
  = 0{,}57721566490153286\ldots
\]
Sie tritt allgegenwärtig in Analysis und Zahlentheorie auf: in der
Laurent-Entwicklung der Riemannschen Zeta-Funktion bei $s = 1$,
\[
  \zeta(s) = \frac{1}{s-1} + \gamma + O(s-1),
\]
im Weierstraßschen Produkt der Gammafunktion,
\[
  \Gamma(z) = \frac{e^{-\gamma z}}{z} \prod_{n=1}^\infty
  \left(1 + \frac{z}{n}\right)^{-1} e^{z/n},
\]
sowie in der Digamma-Funktion: $\psi(1) = \Gamma'(1)/\Gamma(1) = -\gamma$.

Trotz ihrer zentralen Rolle in der Mathematik ist die arithmetische Natur
von $\gamma$ im Wesentlichen unbekannt:

\begin{conjecture}[Irrationalität von $\gamma$]
\label{verm:gamma-irr}
Die Euler-Mascheroni-Konstante $\gamma$ ist irrational.
\end{conjecture}

\begin{conjecture}[Transzendenz von $\gamma$]
\label{verm:gamma-transz}
Die Euler-Mascheroni-Konstante $\gamma$ ist transzendent.
\end{conjecture}

Beide Vermutungen sind Stand 2026 völlig offen.

\begin{bemerkung}
Papanikolaou (1997) zeigte: Falls $\gamma = p/q$ rational mit $\ggT(p,q) = 1$
ist, so gilt $q > 10^{242080}$. Dies ist kein Beweis der Irrationalität --
es schließt lediglich Nenner bis zu dieser Schranke aus. Das Modul
\texttt{euler\_gamma\_irrationality.py} implementiert
\texttt{rational\_approximation\_bound()}, das dieses Resultat zurückgibt.
\end{bemerkung}

Die Verbindung zwischen $\gamma$ und der Zeta-Funktion ist in den
Stieltjes-Konstanten $\gamma_n$ kodiert, definiert durch die Laurent-Entwicklung
\[
  \zeta(s) = \frac{1}{s-1} + \sum_{n=0}^\infty \frac{(-1)^n \gamma_n}{n!}
  (s-1)^n, \qquad \gamma_0 = \gamma.
\]
Die Irrationalität einer jeden Stieltjes-Konstante $\gamma_n$ ist ebenfalls
unbekannt.

Weitere Reihendarstellungen von $\gamma$:
\begin{itemize}[leftmargin=2em]
  \item \textbf{Vacca-Formel (1910):} $\gamma = \sum_{k=1}^\infty (-1)^k
        \lfloor \log_2 k \rfloor / k$
  \item \textbf{Euler/Zeta-Darstellung:} $\gamma = 1 - \sum_{n=2}^\infty
        (\zeta(n) - 1)/n$
  \item \textbf{Brent-McMillan-Algorithmus:} Effiziente Berechnung auf
        $O(n)$ Dezimalstellen in $O(n \log^3 n)$ Zeit.
\end{itemize}


%% ==========================================================================
\section{Offene Probleme}
\label{sec:offen}
%% ==========================================================================

Die Transzendenztheorie kennt viele berühmte offene Probleme. Wir nennen die
wichtigsten:

\begin{conjecture}[Transzendenz von $e + \pi$]
\label{verm:e+pi}
Die Zahl $e + \pi$ ist transzendent (insbesondere irrational).
\end{conjecture}

\begin{conjecture}[Transzendenz von $e \cdot \pi$]
\label{verm:e-pi}
Die Zahl $e \cdot \pi$ ist transzendent.
\end{conjecture}

\begin{bemerkung}
Es ist bekannt, dass $e + \pi$ und $e \cdot \pi$ nicht beide algebraisch sein
können (da sie Summe und Produkt der Nullstellen von $z^2 - (e+\pi)z + e\pi$
sind, und wären beide algebraisch, so auch $e$ und $\pi$ -- Widerspruch).
Also ist mindestens eine von beiden transzendent, aber keine einzelne Aussage
wurde bisher bewiesen.
\end{bemerkung}

\begin{conjecture}[Algebraische Unabhängigkeit von $e$ und $\pi$]
\label{verm:e-pi-unabh}
Die Zahlen $e$ und $\pi$ sind algebraisch unabhängig über $\mathbb{Q}$. Das
heißt, es gibt kein von null verschiedenes Polynom $P(x,y) \in \mathbb{Q}[x,y]$
mit $P(e,\pi) = 0$.
\end{conjecture}

\begin{conjecture}[Transzendenz von $\zeta(3)/\pi^3$]
\label{verm:zeta3}
Die Zahl $\zeta(3)/\pi^3$ ist transzendent. Es ist noch nicht einmal bekannt,
ob $\zeta(3)/\pi^3$ irrational ist (obwohl $\zeta(3)$ selbst irrational ist --
bewiesen von Apéry, 1978).
\end{conjecture}

\begin{conjecture}[Schanuels Vermutung, 1960er Jahre]
\label{verm:schanuel}
Sind $z_1, \ldots, z_n \in \mathbb{C}$ linear unabhängig über $\mathbb{Q}$,
dann haben die $2n$ Zahlen $z_1, \ldots, z_n, e^{z_1}, \ldots, e^{z_n}$
mindestens Transzendenzgrad $n$ über $\mathbb{Q}$.
\end{conjecture}

\begin{bemerkung}
Schanuels Vermutung würde, wenn sie bewiesen wäre, nahezu alle bekannten
Transzendenzergebnisse und weit mehr implizieren:
\begin{itemize}[leftmargin=2em]
  \item $e$ und $\pi$ sind algebraisch unabhängig (mit $z_1 = 1$, $z_2 = i\pi$).
  \item Der Satz von Lindemann--Weierstrass folgt direkt.
  \item Die Transzendenz von $e^\pi$, $\ln\pi$, $\ln\ln 2$ usw.\ würde folgen.
\end{itemize}
Trotz allgemeiner Überzeugung von ihrer Richtigkeit wurde Schanuels Vermutung
bisher nicht bewiesen -- nicht einmal in Spezialfällen über das hinaus, was
Bakers Satz liefert.
\end{bemerkung}

Tabelle~\ref{tab:status} fasst den Transzendenzierstatus wichtiger Konstanten
zusammen.

\begin{table}[ht]
\centering
\caption{Transzendenzierstatus ausgewählter Konstanten (Stand 2026)}
\label{tab:status}
\begin{tabular}{lll}
\toprule
Konstante & Status & Quelle \\
\midrule
$e$ & Transzendent (bewiesen) & Hermite (1873) \\
$\pi$ & Transzendent (bewiesen) & Lindemann (1882) \\
$e^\pi$ & Transzendent (bewiesen) & Gelfond--Schneider \\
$2^{\sqrt{2}}$ & Transzendent (bewiesen) & Gelfond--Schneider \\
$\ln 2$, $\ln 3$ & Transzendent (bewiesen) & Lindemann--Weierstrass \\
$e + \pi$ & Offen (vermutet transzendent) & --- \\
$e \cdot \pi$ & Offen (vermutet transzendent) & --- \\
$\gamma$ & Irrationalität offen & Papanikolaou (1997) \\
$\zeta(3)$ & Irrational (bewiesen), Transzendenz offen & Apéry (1978) \\
$\zeta(5), \zeta(7), \ldots$ & Offen & --- \\
\bottomrule
\end{tabular}
\end{table}


%% ==========================================================================
\section{Algorithmische Aspekte und Implementierung}
\label{sec:algorithmisch}
%% ==========================================================================

Das Modul \texttt{transcendence\_theory.py} (Build 141, Specialist Mathematics
Project) stellt numerische Werkzeuge zur Untersuchung von Transzendenzeigenschaften
bereit.

\paragraph{Zahlenklassifikation.}
Die Funktion \texttt{classify\_number(x)} versucht zu entscheiden, ob $x$
rational, algebraisch oder wahrscheinlich transzendent ist. Für einen
Gleitkommawert prüft sie zunächst Rationalität via bester rationaler
Approximation (\texttt{fractions.Fraction}), sucht dann nach einem
Minimalpolynom vom Grad $\leq 6$ mit ganzzahligen Koeffizienten in
$\{-8, \ldots, 8\}$.

\paragraph{Irrationalitätsmaß.}
Die Funktion \texttt{irrationality\_measure(x)} schätzt $\mu(x)$, indem sie
$\mu \approx -\log|x - p/q|/\log q$ für die besten rationalen Näherungen
$p/q$ mit $q \leq 10000$ berechnet.

\paragraph{Kettenbruchanalyse.}
Die Funktion \texttt{continued\_fraction\_irrationality(x)} berechnet die
Kettenbruchentwicklung $[a_0; a_1, a_2, \ldots]$ und schätzt das
Irrationalitätsmaß aus dem Wachstum der Teilnenner $a_k$ anhand der Beziehung
$\mu \approx 2 + \limsup_k \log(a_{k+1})/\log(q_k)$.

\paragraph{Euler-Mascheroni-Konstante.}
Die Klasse \texttt{EulerMascheroniConstant} (in
\texttt{euler\_gamma\_irrationality.py}) berechnet $\gamma$ mit beliebiger
Präzision via \texttt{mpmath} (Brent-McMillan-Algorithmus), verifiziert die
Papanikolaou-Schranke und illustriert die Verbindung zur Gammafunktion
über $\psi(1) = \Gamma'(1) = -\gamma$.


%% ==========================================================================
\section{Literatur}
%% ==========================================================================

\begin{thebibliography}{99}

\bibitem{hermite1873}
C.~Hermite,
\textit{Sur la fonction exponentielle},
Comptes Rendus de l'Acad\'emie des Sciences \textbf{77} (1873), 18--24, 74--79,
226--233, 285--293.

\bibitem{lindemann1882}
F.~von Lindemann,
\textit{\"Uber die Zahl $\pi$},
Mathematische Annalen \textbf{20} (1882), 213--225.

\bibitem{weierstrass1885}
K.~Weierstrass,
\textit{Zu Lindemann's Abhandlung: ``\"Uber die Ludolph'sche Zahl''},
Sitzungsberichte der K\"oniglich Preu\ss{}ischen Akademie der Wissenschaften
\textbf{5} (1885), 1067--1086.

\bibitem{gelfond1934}
A.O.~Gelfond,
\textit{Sur le septi\`eme probl\`eme de Hilbert},
Bulletin de l'Acad\'emie des Sciences de l'URSS, S\'erie Math\'ematique
\textbf{7} (1934), 623--630.

\bibitem{schneider1935}
T.~Schneider,
\textit{Transzendenzuntersuchungen periodischer Funktionen, I},
Journal f\"ur die reine und angewandte Mathematik \textbf{172} (1935), 65--69.

\bibitem{baker1966}
A.~Baker,
\textit{Linear forms in the logarithms of algebraic numbers},
Mathematika \textbf{13} (1966), 204--216.

\bibitem{baker1975}
A.~Baker,
\textit{Transcendental Number Theory},
Cambridge University Press, Cambridge, 1975.

\bibitem{papanikolaou1997}
T.~Papanikolaou,
\textit{Einige neue Resultate \"uber den Euler-Mascheroni-Quotienten},
Diplomarbeit, Universit\"at des Saarlandes, 1997.

\bibitem{waldschmidt2000}
M.~Waldschmidt,
\textit{Diophantine Approximation on Linear Algebraic Groups},
Grundlehren der Mathematischen Wissenschaften \textbf{326},
Springer, Berlin, 2000.

\bibitem{nesterenko1996}
Yu.\,V.~Nesterenko,
\textit{Modular functions and transcendence},
Sbornik: Mathematics \textbf{187} (1996), 1319--1348.

\end{thebibliography}

\end{document}
